\documentclass[10pt]{article}
\usepackage[margin=0.68in]{geometry}
\usepackage{amsmath,amssymb}
\usepackage[T1]{fontenc}
\usepackage[hidelinks]{hyperref}

\begin{document}
\begin{center}
{\Large Literature Status: Finite Multi-Utilities Need Not Admit\\Finite Injective Refinements\par}
\vspace{0.6em}
{fa\_banach\_001, lane 8\par}
\vspace{0.25em}
{August 9, 2026\par}
\end{center}
\vspace{0.7em}

\paragraph{Status.}
\textbf{literature\_already\_answered}. The later paper supplies a stronger negative answer.

\paragraph{Original source and question.}
Hack, Braun, and Gottwald, \emph{The classification of preordered spaces in terms of monotones: complexity and optimization}, arXiv:2202.12106; Theory and Decision 94 (2023), 103--125. A family $V$ is a multi-utility for a preorder $\preceq$ when
\[
 x\preceq y \quad\Longleftrightarrow\quad v(x)\leq v(y)\ \text{for every }v\in V.
\]
A member $v$ is an injective monotone when it is monotone and $v(x)=v(y)$ implies $x\sim y$. After Proposition 10 (labelled \texttt{finite relations} in the arXiv source), and again in the final ``Open questions'' paragraph, the paper asks whether every preorder with a finite multi-utility has a finite injective-monotone multi-utility, or whether a counterexample exists.

\paragraph{Supporting answer.}
Hack, Braun, and Gottwald, \emph{The infinite information gap between mathematical and physical representations}, arXiv:2203.16272v4 (earlier title: \emph{On a geometrical notion of dimension for partially ordered sets}). The paper calls multi-utilities ``mathematical representations,'' strict-monotone multi-utilities ``physical representations,'' and injective-monotone multi-utilities ``injective physical representations.'' Theorem 2 constructs a partial order whose minimal multi-utility has two functions but whose minimal strict-monotone multi-utility is countably infinite. Because every injective monotone is a strict monotone, this already rules out every finite injective-monotone multi-utility.

The same conclusion is stated directly in the paper's dimension language. Lemma 1 identifies the minimum cardinality of an injective-monotone multi-utility with Debreu dimension. Theorem 7 then states that there are partial orders with finite geometrical dimension (the minimum size of a multi-utility) and countably infinite Debreu dimension. Hence the source question has a full negative answer.

\paragraph{Identification and scope.}
This is a separate later arXiv paper by the same authors. It cites the classification work and states that these results improve the classification of preordered spaces. The answer is full for finite multi-utilities versus finite injective-monotone multi-utilities. The later paper separately leaves open how finite strict-monotone and finite injective-monotone multi-utilities compare; that narrower question is not claimed resolved here.

\paragraph{Evidence searched.}
The run indexes contained no prior record for the source question. Searches combining finite multi-utilities, injective monotones, and Debreu dimension located arXiv:2203.16272. The decisive locations are Theorem 2, Lemma 1, Theorem 7, and Appendix C.

\begin{thebibliography}{9}
\bibitem{HBGclass}
P. Hack, D. A. Braun, and S. Gottwald,
\emph{The classification of preordered spaces in terms of monotones: complexity and optimization},
arXiv:2202.12106; Theory and Decision 94 (2023), 103--125.

\bibitem{HBGgap}
P. Hack, D. A. Braun, and S. Gottwald,
\emph{The infinite information gap between mathematical and physical representations},
arXiv:2203.16272v4.
\end{thebibliography}

\end{document}
